\chapter{Wstęp}
\section{Wprowadzenie}
% TO DO: poprawiłem nieco styl wypowiedzi w poniższych akapitach. Ale nie zrobiłem wszystkiego. Proszę uważniej formułować zdania, by nikt nie mógł Panu zarzucić brak poprawności gramatycznej.

Osoby podróżujące do miejsca pracy lub na uczelnię niemal codziennie muszą mierzyć się z problemem dotarcia do celu. Jeśli cel jest blisko, zazwyczaj docierają do niego pieszo. Jednak w przypadku dużych odległości czy też jakichś niedogodności muszą  korzystać ze jakichś środków transportu.

% TO DO: proszę odpowiednio oznaczać kropki stojące za skrótami by odróżniać je od kropek kończących zdanie (zamieszcając ukośnik i spację zaraz za kropką albo tyldę łączącą się z następnym wyrazem). Pisałem o tym w szablonie. 

Dostępną wtedy opcją jest komunikacja miejska, np.\ tramwaje lub autobusy. Jest ona tania, a częstotliwość kursów i dostępność w obrębie centrum miast jest zadowalająca. Wadami tego rozwiązania są częste awarie tego transportu, wykolejenia, wypadki, które potrafią sparaliżować trasę wielu linii komunikacyjnych oraz niska dostępność w miejscach oddalonych od miasta i wsiach. W tych miejscach linie potrafią kursować bardzo sporadycznie, co przyczynia się do tzw.\ wykluczenia komunikacyjnego, które jest szczególnie odczuwalne wśród osób starszych, niepełnosprawnych i dzieci.

% TO DO: proszę przed wrzuceniem zmian sprawdzić literówki (jakimś spell-checkerem, ja korzystam z wbudowanego spell-checkera w texniccenter)
Alternatywą jest dojazd samochodem. Daje on większą wygodę i pozwala ominąć problem niskiej częstotliwości kursów komunikacji miejskiej. Co więcej, dzięki nowoczesnym systemom nawigacyjnym planującym trasę można nim osiągnąć cel nawet przy wystąpieniu utrudnień na drodze. 

Koszty związane z tą formą przejazdów są zdecydowanie wyższe niż w przypadku komunikacji miejskiej. Składają się na nie opłaty za paliwo, ubezpieczenie, przeglądy, parkingi oraz koszty serwisów. Samo prowadzenie samochodu nie jest dostępne dla osób bez prawa jazdy, które nie mogą go uzyskać ze względów zdrowotnych lub finansowych. Dużym problemem są również emitowane przez samochody spaliny, które przyczyniają się do zanieczyszczenia powietrza i globalnego ocieplenia.

Istnieją również inne metody transportu, takie jak rower, albo skuter, jednak czynniki atmosferyczne, takie jak deszcz, śnieg czy mróz potrafią skutecznie uniemożliwić korzystania z tych środków transportu.

Pewną receptą na wysokie koszty dojazdu samochodem oraz niską dostępność komunikacji miejskiej jest tzw.\ ,,współużytkowanie pojazdów'' (ang.~\emph{carpooling}). Polega ono na wzajemnym podwożeniu czy umawianiu się na wspólne przejazdy z ustalonymi zasadami ich rozliczania. Ta forma współpracy pozwala zapewnić wygodę, jak również szybkość podróży.
Czasem jednak szwankuje z uwagi na niedostępność kierowców lub pasażerów chcących podróżować w tym samym kierunku.

% TO DO: proszę dołożyć parę zdań o tym, że sam pomysł na współdzielenie zasobów nie jest czymś nowym (rowery miejskie i hulajnogi można pod to podciągnąć). Zapewne istnieją gotowe aplikacje pozwalające "umawiać się" na przejazdy. Proszę wskazać te aplikacje. Proszę dołożyć jakieś cytowania.
% Powinien Pan napisać, dlaczego zamiast się nimi posłużyć zdecydował Pan opracować własną aplikację. Powinien Pan przedstawić motywację, jaka za tym stała.


% TO DO: proszę przed dostarczeniem zmian do repozytorium sprawdzić literówki. Ja korzystam z texniccenter, który ma podpięty polski słownik i pomaga eliminować takie błędy. Jest to prosty spell-checker, ale daje radę.

\section{Cel i zakres pracy}
Celem pracy jest zaprojektowanie oraz zaimplementowanie aplikacji internetowej wspierającej użytkowników w organizowaniu wspólnych przejazdów.
Aplikacja ta powinna umożliwiać użytkownikom dopasowanie pasażerów do kierowców, jak również pozwalać na wyznaczanie optymalnych tras, 
uwzględniając preferencje jak i początkowe oraz docelowe lokalizacje wszystkich uczestników przejazdu. 
Poza rozwiązaniem problemu wykluczenia komunikacyjnego wdrożenie aplikacji powinno dać szansę na wykształcenia nawyków pro-ekologicznych, 
wynikających ze współdzielenia zasobów, przekładających się na zmniejszenie śladu węglowego oraz jednostkowego kosztu przejazdu.

W zakres pracy wchodzi zaprojektowanie, zaimplementowanie, przetestowanie oraz wdrożenie bazy danych oraz REST API, które będzie obsługiwało aplikację internetową. 
% TO DO: czy API obsługuje aplikację, czy raczej aplikacja korzysta z API? Niby mowa o tym samym, jednak dalej Pan pisze, że to aplikacja korzysta z "zaimpelemtowanych funkcjonalności". - proszę doczytać zalecenia w szablonie, pkt 6 na str. 29.

W skład API wchodzić będzie logika biznesowa, warstwa zabezpieczeń dostępu do punktów końcowych, jak również warstwa dostępu do bazy danych. 
W ramach pracy zostanie również zaprojektowana i zaimplementowana aplikacja internetowa, która pozwoli na korzystanie z zaimplementowanych funkcjonalności.

% Czy algorytm oparty o A* jest heurystyką, czy metaheurystyką? Wydaje mi się, że jest to heurystyka, ale nie jestem pewien. 
% W przebieg algorytmu wchodzi wybranie najbliższego kierowcy, którego punkty początkowe i końcowe są najbliższe punktom początkowym i końcowym pasażera.
Proces dobierania pasażerów do kierowców będzie polegał na metaheurystyce/heurystyce opartej o algorytm "A-Star", który będzie również wyznaczał optymalną trasę przejazdu.


\section{Układ pracy}
